\section{Chern intersections and boundary limits of a highest Hodge line}
\label{p1varhd:section}

We prove the Chern-intersection formula deferred in
Lemma~\ref{p1logalt:rank}, and then give a different proof of the
boundary restriction estimate in Lemma~\ref{p1logalt:boundary-drop}.  Its two ingredients are a Chern-intersection
formula with arbitrary smooth closed test forms and a family of Thom
forms concentrating on a boundary component.  The limiting Hodge norm
computes the tangential curvature at a general boundary point; a product
Poincar\'e bound controls the crossings and permits passage to the
restriction intersection number.

Throughout this section, $R$ is a smooth connected projective complex
variety, $D_R$ is a reduced simple-normal-crossing divisor, and
$\mathbb V$ is a polarizable integral pure variation on $R\setminus D_R$.
A highest Hodge line is its highest nonzero filtration piece $F^p$, with
$F^{p+1}=0$ and $\rank F^p=1$.  We denote its rational parabolic
extension by $A$.  The parabolic convention includes finite-monodromy
weights.  Its line period map is the map from the universal cover of
$R\setminus D_R$ to the projective space of a fixed flat reference fiber
which remembers $F^p$.  Its rank need not equal the full-period rank.
For $n=\dim R$ we use
\[
 \nu(A)=\max\{j\in\{0,\ldots,n\}:
                   c_1(A)^j\cdot H^{n-j}>0\},
\]
where $H$ is any ample class and $A$ is nef.  In particular, this
numerical dimension is zero on a point.  Normalize
$\ddc=\frac{i}{2\pi}\partial\bar\partial$.

Lemma~\ref{p1logalt:extensions} supplies nefness and compatibility of
the rational extension with positive tensor powers and log-pair
pullbacks whose interiors map into the original interior. The
Chern-intersection formula below uses these extension properties,
but neither full-period descent nor a boundary rank-drop result.
For strict rank loss we also use the integral tensor replacement
of Lemma~\ref{p1logalt:kernel} and the transverse-loop criterion of
Lemma~\ref{p4int:transverse-loop}. The replacement is faithful on
the moving highest line: a monodromy element fixing that line at
every point of the connected universal cover is the identity.
The Thom-form calculation computes restriction intersections from
tangential limiting curvature. It uses the local graded Hodge data,
without identifying the boundary restriction with a globally extended
graded line as in Section~\ref{p1logalt:section}.

\subsection{Curvature and intersection powers}

\begin{lemma}[Chern intersections with closed test forms]\label{p1varhd:curvature}
Let $A$ be the rational parabolic highest line of an integral
polarized pure variation on the complement of a reduced
simple-normal-crossing divisor in a smooth projective variety $R$.
Then
\begin{equation}\label{p1varhd:numrank}
 \nu(A)=\text{generic differential rank of its projective line map}
\end{equation}
provided $A$ is nef.  In the unipotent case, if $\alpha$ is the Chern
form of its Hodge metric on $R\setminus D_R$, then for every
$0\leq j\leq n=\dim R$ and every smooth closed form $\gamma$ of type
$(n-j,n-j)$, the integral
$\int_{R\setminus D_R}\alpha^j\wedge\gamma$ is absolutely convergent
and equals $c_1(A)^j\cdot[\gamma]$.
\end{lemma}

\begin{proof}
We may compute after a generically finite smooth alteration on
which all boundary monodromies are unipotent.  To obtain one, pass
to a finite level cover with neat monodromy, compactify, and resolve.
Quasi-unipotence then gives unipotence, including for the commuting
local monodromies on the resolved boundary.  The parabolic pullback
property of Lemma~\ref{p1logalt:extensions} identifies the new line
with the pullback of $A$.
Generically finite pullback preserves nef numerical dimension by
the intersection and projection formulas.  It also preserves the
generic differential rank in characteristic zero.  We therefore
work on the unipotent model, retaining the notation $R,A$.

Let $\alpha$ be the semipositive Chern form of the highest-line
Hodge metric on the open part. Griffiths' curvature formula
\cite[Theorem~5.2, Formula~5.3, and Formula~6.18]{GriffithsPeriodsIII} says
that $\alpha(v,\bar v)$ is a positive fixed multiple of the squared
norm of the projective highest-line differential in direction $v$.
In particular, its rank is the rank in \eqref{p1varhd:numrank}.

Consider a crossing chart with coordinates $z_1,\ldots,z_n$ and
boundary $z_1\cdots z_c=0$.  Put
$\rho_i=-\log|z_i|^2$ and shrink the chart so that $\rho_i>e$.
The horizontal Schwarz estimate, applied separately to the
coordinate disks, gives
\begin{equation}\label{p1varhd:poincare}
 0\leq\alpha\leq C\Omega,
 \qquad
 \Omega=
 \sum_{i=1}^{c}\frac{i\,dz_i\wedge d\bar z_i}
                         {|z_i|^2\rho_i^2}
 +\sum_{i=c+1}^{n}i\,dz_i\wedge d\bar z_i.
\end{equation}
Here and below harmless positive normalization constants are
absorbed into $C$.  The coordinate estimate follows from the
curvature of the horizontal period metric and the singular
Ahlfors--Schwarz lemma; see \cite[\S 2.1, Proposition~2.4 and
Lemma~2.5]{BrunebarbeCadorel}.  Positivity and the Hermitian-form Cauchy--Schwarz
inequality control the off-diagonal entries and give the uniform
product estimate \eqref{p1varhd:poincare}, even where the period
differential is degenerate.

By the nilpotent-orbit theorem
\cite[Theorem~4.12]{Schmid}, a nonvanishing top extension frame,
written in flat coordinates on a bounded angular sector, is
\begin{equation}\label{p1varhd:frame}
 v(z)=\exp\left(\sum_{i=1}^{c}
               \frac{\log z_i}{2\pi i}N_i\right)u(z),
\end{equation}
where the $N_i$ commute and $u$ is holomorphic and nonvanishing.
If $\beta$ is the curvature of a fixed smooth metric on $A$, write
$\alpha-\beta=\ddc\varphi$.  The ratio of the two metrics satisfies
\begin{equation}\label{p1varhd:loglog}
 |\varphi(z)|\leq C\left(1+\sum_{i=1}^{c}\log\rho_i\right).
\end{equation}
This is the logarithmic growth estimate for canonical Hodge frames;
see \cite[Theorem~5.1]{CattaniKaplanDegenerating}.  Here the Schwarz
estimate also gives a direct proof of the bound needed for integration.  The differential of the logarithmic Hodge norm of any
flat vector is bounded by a constant times the invariant period
metric, with a constant independent of the point by homogeneity.
Travel from a fixed point along coordinate paths in a bounded
angular sector.  Their product-Poincar\'e length is bounded by
$C(1+\sum_i\log\rho_i)$, so the Hodge metric and its inverse in a
flat frame grow at most as products of powers of the $\rho_i$.
The exponential in \eqref{p1varhd:frame} and its inverse are polynomials
in the logarithms, since the $N_i$ are commuting nilpotent
operators.  Finally, $u$ has ordinary norm bounded above and below
on a smaller chart.  These three observations give the upper and
lower logarithmic bounds in \eqref{p1varhd:loglog}.  We have therefore
obtained the two controls needed for the cutoff argument:
\begin{equation}\label{p1logalt:metric-bounds}
 0\leq\alpha\leq C\Omega,
 \qquad |\varphi|\leq C\Bigl(1+\sum_i\log\rho_i\Bigr).
\end{equation}

Chern-current representation for canonical Hodge extensions is
established in \cite[Theorem~5.2]{CattaniKaplanDegenerating}. We now
justify the highest-line intersection formula with closed test forms,
including the absence of boundary contributions.  For $0\leq j\leq n$ and every
smooth closed form
$\gamma$ of type $(n-j,n-j)$,
\begin{equation}\label{p1varhd:chern}
 \int_{R\setminus D_R}\alpha^j\wedge\gamma
       =c_1(A)^j\cdot[\gamma].
\end{equation}
The case $j=0$ is immediate, so assume $j>0$.  Choose a smooth
function $\chi$ that is one on $(-\infty,1]$ and
zero on $[2,\infty)$, and use products of
$\chi((\log\rho_i)/M)$ to cut off the boundary.  The $\rho_i$ can
be globalized using smooth metrics on the boundary divisor
bundles.  Denote the resulting cutoff by $\chi_M$.  Its second
derivatives are uniformly bounded in the product-Poincar\'e metric,
and their support escapes to the boundary as $M\to\infty$.

Since all curvature forms are closed, on the open part we have
\[
 \alpha^j-\beta^j
 =\ddc\left(\varphi
       \sum_{k=0}^{j-1}\alpha^k\wedge\beta^{j-1-k}\right).
\]
Multiplying by $\chi_M\gamma$ and integrating by parts makes the
difference of the integrals equal to
\[
 \int\varphi\,\ddc\chi_M\wedge
       \sum_{k=0}^{j-1}\alpha^k\wedge\beta^{j-1-k}\wedge\gamma.
\]
Its absolute value is bounded on escaping tails by a constant times
\[
 \left(1+\sum_i\log\rho_i\right)\Omega^n.
\]
This is integrable: in a cusp variable the relevant radial
integrals are bounded by
$\int^{\infty}(1+\log\rho)\rho^{-2}\,d\rho$.
Thus the error tends to zero.  The same domination permits removal
of $\chi_M$ from the two curvature integrals, proving
\eqref{p1varhd:chern}.  Rational bundles are handled by first taking an
integral tensor power and then dividing the forms and identities
by that power.

Let $r$ be the generic rank of the projective line map.  Then
$\alpha^{r+1}=0$ on the open part, whereas $\alpha^r$ is nonzero
and positive on a nonempty open subset.  Apply \eqref{p1varhd:chern}
with complementary powers of an ample K\"ahler form.  It gives
vanishing of the $(r+1)$st intersection power and strict positivity
of the $r$th one.  The nef intersection characterization proves
\eqref{p1varhd:numrank}.
\end{proof}

\subsection{Rank loss at a monodromy divisor}

\begin{lemma}[Boundary rank drop by a Thom-form calculation]\label{p1varhd:boundary}
In the setting of Lemma~\ref{p1varhd:curvature}, suppose that monodromy
is faithful on the moving highest line.  If $D_i$ is a boundary
component with nonidentity transverse monodromy, then
\begin{equation}\label{p1varhd:strictdrop}
             \nu(A|_{D_i})<\nu(A).
\end{equation}
\end{lemma}

\begin{proof}
Put $r=\nu(A)$.  We first compute a bound for the numerical
dimension on a boundary component of a unipotent model.  This
computation is unchanged on a component dominating $D_i$ under a
generically finite alteration, since restriction of the pulled-back
line is the generically finite pullback of $A|_{D_i}$.

\paragraph{The limiting highest line.}
Near a general point of a unipotent boundary component write the
coordinates as $(t,s)$, with divisor $t=0$ and monodromy logarithm
$N$.  Let $w$ be the weight and $p$ the largest filtration index.
The limiting mixed Hodge structure is polarized by $N$
\cite[Theorem~6.16]{Schmid}; its graded structures vary as polarized
variations along the smooth boundary stratum
\cite[Proposition~2.10 and the following discussion]{CattaniKaplanDegenerating}.
The limiting $F^p$ has dimension one and $F^{p+1}=0$.  Consequently
it consists of one Deligne summand $I^{p,q}$.  Write
$p+q=w+l$.  In the primitive decomposition for $N$, a nonzero
$N^k$-image term with $k>0$ would have to come from $F^{p+k}$,
which is zero.  The line is therefore primitive, $l\geq0$, and,
for a nonzero vector $u_0(s)=u(0,s)$ spanning it,
\[
       N^lu_0(s)\ne0,\qquad N^{l+1}u_0(s)=0.
\]
Shrink the boundary stratum so that its Hodge type, and hence $l$,
is fixed.

For a highest Hodge vector in the pure variation, the Hodge norm is
the flat polarization pairing with its conjugate, multiplied by the
fixed factor determined by its Hodge type.
Substituting the exponential frame into that pairing gives a
polynomial in $y=-\log|t|$.  Primitive polarization identifies its
leading term as a positive constant times $y^l$ times the squared
norm $h_l(s)$ of the primitive graded line.  Replacing $u(t,s)$
by $u_0(s)$ introduces errors $O(|t|)$ times powers of $y$;
the same holds after tangential differentiation, by holomorphicity
of $u$.  It follows, locally uniformly with tangential derivatives,
that
\begin{equation}\label{p1varhd:limitnorm}
 \log\|v(t,s)\|^2
      =l\log y+\log h_l(s)+c+O(y^{-1})+O(|t|y^M)
\end{equation}
for some $M$ and a constant $c$.  Thus the tangential curvature
converges to the curvature of the primitive graded highest line.
If $q_i$ denotes the generic rank of its projective differential,
the limiting tangential curvature has rank at most $q_i$.

\paragraph{The restriction intersection and crossing points.}
We claim that
\begin{equation}\label{p1varhd:restrictionrank}
                    \nu(A|_{D_i})\leq q_i.
\end{equation}
The following estimate proves the claim without assuming uniform
convergence of \eqref{p1varhd:limitnorm} at crossings.  Fix an ample
K\"ahler form $\omega$ and an integer $j$ with
$q_i<j\leq n-1$.  Choose a smooth metric on $\cO(D_i)$, write its
curvature as $\beta_i$, and put
$\tau=\log\|s_{D_i}\|^2$.  Let $g$ be a fixed smooth convex
function with $g'=0$ on $(-\infty,0]$ and $g'=1$ on
$[\log4,\infty)$, and set
$f_\epsilon(x)=g(x-\log\epsilon^2)$.  Then
\begin{equation}\label{p1varhd:thom}
 \eta_\epsilon
  =\beta_i+\ddc f_\epsilon(\tau)
  =(1-f_\epsilon'(\tau))\beta_i
     +\frac{i}{2\pi}f_\epsilon''(\tau)
                   \partial\tau\wedge\bar\partial\tau
\end{equation}
is a smooth closed representative of $c_1(\cO(D_i))$.
Here the normalization is $\ddc\log|t|^2=[t=0]$.
The function $f_\epsilon(\tau)$ is constant near $D_i$, so it
extends smoothly there.  Outside $\|s_{D_i}\|\leq2\epsilon$,
the two curvature terms cancel and $\eta_\epsilon=0$.

On a crossing chart put
\[
 \lambda=\frac{i\,dt\wedge d\bar t}{|t|^2},\qquad
 \rho=-\log|t|^2,\qquad
 \Omega=\frac{\lambda}{\rho^2}+\Omega_s,
 \qquad T=\alpha^j\wedge\omega^{n-1-j}.
\]
The tangential form $\Omega_s$ is a product cusp metric, with
Euclidean summands in the remaining coordinates.  It has finite
volume.  The bound \eqref{p1varhd:poincare} gives
$0\leq T\leq C\Omega^{n-1}$.
Write $\tau=\log|t|^2+\psi$ with $\psi$ smooth, and put
$L_\epsilon=|\log\epsilon|$.  The smooth first term of
\eqref{p1varhd:thom} contributes at most
\[
 C\int_{|t|\leq C\epsilon}\Omega^n
                     =O(L_\epsilon^{-1}).
\]
For the second term write $\partial\tau=a_0+b_0$, with
$a_0=dt/t$ and $b_0=\partial\psi$.  Its support lies in
$c\epsilon\leq|t|\leq C\epsilon$, with constants independent
of $s$.  On this shell, positivity gives
\begin{align*}
 \int i a_0\wedge\bar a_0\wedge T&=O(1),\\
 \int i b_0\wedge\bar b_0\wedge T
       &\leq C\int_{c\epsilon\leq|t|\leq C\epsilon}\Omega^n
         =O(L_\epsilon^{-2}).
\end{align*}
The first estimate follows because wedging with $\lambda$
selects only the tangential part of $T$.  For the second, a smooth
positive form dominating $i b_0\wedge\bar b_0$ is bounded by
$C\Omega$.  Cauchy--Schwarz for the positive form $T$ now bounds
the mixed terms by $O(L_\epsilon^{-1})$.
The bounded coefficient $f_\epsilon''$ does not change these
estimates.

It remains to consider the pure normal-shell term.  Set
$t=\epsilon u$.  Its cutoff coefficient becomes
\[
 g''\bigl(\log|u|^2+\psi(\epsilon u,s)\bigr),
\]
whose support is contained in a fixed annulus, since $\psi$ is
bounded.  The absolute integrand is bounded, independently of
$\epsilon$, by
\[
 C\,\frac{i\,du\wedge d\bar u}{|u|^2}
            \wedge\Omega_s^{n-1}.
\]
This is integrable on that annulus times the tangential chart,
including all tangential divisor crossings.  For almost every $s$,
\eqref{p1varhd:limitnorm} gives a tangential curvature limit of rank
less than $j$, so the tangential $j$th power tends to zero.
Dominated convergence makes the pure normal-shell integral tend
to zero.  The estimates are uniform under truncating the
tangential cusp variables, so no mass is lost when those
truncations are removed.  Choose a finite partition of unity subordinate to crossing charts.
Multiplication by its bounded smooth functions preserves each absolute
estimate.  Summing the local estimates proves the same conclusion on
all of $D_i$; no integration by parts is performed after inserting this
partition.

For each fixed $\epsilon$, apply \eqref{p1varhd:chern} with
$\gamma=\omega^{n-1-j}\wedge\eta_\epsilon$.  Its cohomology
class is independent of $\epsilon$.  The preceding limit gives
\[
 c_1(A|_{D_i})^j\cdot[\omega|_{D_i}]^{n-1-j}=0.
\]
Nefness proves \eqref{p1varhd:restrictionrank}.  If $q_i=n-1$, that
inequality is automatic.  This also explains why the generic
limit calculation alone would not be sufficient: the uniform
product bound is what controls the crossing locus.

\paragraph{A fixed quotient and transverse loops.}
Return to the original boundary component, with transverse
monodromy $T_s\exp N$, where $T_s$ has finite order.  Work first
in a transverse root coordinate killing $T_s$.  Set
$C_N=\operatorname{coker}N$, endowed with the quotient of the
fixed monodromy weight filtration.  The primitive limiting line
of weight $k=w+l$ projects nontrivially into $W_kC_N$.
The fixed linear quotient
\begin{equation}\label{p1varhd:fixedquotient}
              W_kC_N\longrightarrow\operatorname{Gr}^W_k C_N
\end{equation}
recovers its primitive graded line.  Indeed, in the Lefschetz
decomposition of $\operatorname{Gr}^W_k$, quotienting by the
image of $N$ removes exactly the summands that are positive powers
of $N$ applied to higher primitives.  The induced quotient is the
primitive part of weight $k$.  In particular, projective
differential rank cannot increase under \eqref{p1varhd:fixedquotient}.
No choice of a moving primitive splitting is involved.

The limiting line in $C_N$ lies in one $T_s$-eigenspace.  Compose
the quotient by $N$ with the fixed projection onto this eigenspace.
Near the chosen general boundary point this gives a nonzero
holomorphic projective map from the untwisted highest line.
Under deck rotation the untwisted line is acted on by $T_s$;
on the projected eigenspace that action is scalar.  The projective
map therefore descends to the original chart.  On the punctured
chart it factors through the moving highest-line map, since the
exponential of $N$ disappears in $C_N$.  Its boundary differential
rank is at least the rank $q_i$ of the primitive graded line.

Suppose $\nu(A|_{D_i})\geq r$.  By
\eqref{p1varhd:restrictionrank}, the projected map has at least $r$
independent tangential differentials at a general boundary point.
Choose $r$ affine components $g_1,\ldots,g_r$ with that property.
Their differentials are independent along the boundary at that point,
and the map $g=(g_1,\ldots,g_r)$ extends holomorphically over the original
chart.  On the punctured chart it factors through the highest-line map.
Lemma~\ref{p1varhd:curvature} gives the latter rank at most $r$.
The hypotheses of the transverse-loop criterion
(Lemma~\ref{p4int:transverse-loop}) are therefore satisfied in the
original transverse coordinate, before the finite root cover.  It
follows that the original transverse monodromy fixes the moving line
on the connected universal cover.  This contradicts faithfulness and
its nonidentity.  When $r=0$, use the same lemma with an empty list of
functions; when $r>n-1$, the supposed inequality
$\nu(A|_{D_i})\geq r$ is impossible on dimensional grounds.  This proves
Equation~\eqref{p1varhd:strictdrop}.
\end{proof}


\subsection{Connection with the integral projective argument}

The hypotheses required for the preceding calculation arise from the
integral tensor replacement and full-period descent in
Lemmas~\ref{p1logalt:kernel} and~\ref{p1logalt:descent}.
On the resulting projective base, Lemma~\ref{p1varhd:boundary} gives the
same strict inequality as Lemma~\ref{p1logalt:boundary-drop}, with
the boundary intersection justified directly by Thom forms.  The
boundary subtraction estimate of Lemma~\ref{p1logalt:volume} then
applies to the union $J$ of the divisors with nonidentity transverse
monodromy.  The full-period map supplies the bigness of $K_R+J$ as in
the proof of Theorem~\ref{p1logalt:positivity}; the highest-line
rank supplies the strict decrease in the degree of the restriction
polynomials.  Thus the two different period ranks retain their separate
roles in adjoint positivity.

\begin{remark}[Tensor replacement on a restricted locus]
\label{p4int:restricted-tensor}
The same local results also give the curve-constancy part
of the highest-line properties used in such comparisons.  On a smooth
complete curve, degree zero gives zero integral of the semipositive
Hodge curvature by Equation~\eqref{p1varhd:chern}, after an adapted
cover if necessary.  The curvature is therefore zero, and the highest
line is projectively constant in a flat reference.  If the highest line
is projectively constant on a connected smooth algebraic locus, apply
Lemma~\ref{p1logalt:kernel} to the restricted variation itself.
Use its own monodromy image and isotypic decomposition: restricting
an ambient replacement need not preserve faithfulness, because the
restricted line image may have a larger projective stabilizer.
Every monodromy element of the replacement fixes its constant highest
line, so faithfulness makes that monodromy trivial.  Since the new line
is a positive tensor power of the original line, the original
highest-line character on that locus has finite image.  The tensor
exponent here is allowed to depend on the restricted variation.
\end{remark}

\begin{remark}\label{p1varhd:further-models}
The adjoint conclusion is preserved under further smooth projective
birational modification $\pi:R'\to R$, together with a compatible
modification of the source.  The parabolic pullback property preserves
the line-bundle identity and nefness, while
\[
 K_{R'}+b\pi^*A
   =\pi^*(K_R+bA)+E_\pi,\qquad E_\pi\geq0,
\]
preserves bigness.  The calculations here concern the whole rank-one
highest step of a polarizable integral pure variation on a projective
pair; they assert no equality between highest-line rank, full-period
rank, and whole-fiber birational variation.
\end{remark}
